% Construct Validity Failure in Single-Cell Embedding Evaluation — poster v1
% tikzposter 3x3 block format, matching grokking_poster_v2.tex style.
%
%   lualatex scib_poster_v1.tex

\documentclass[25pt, a0paper, portrait]{tikzposter}

\usepackage{amsmath,amssymb}
\usepackage{booktabs}
\usepackage{graphicx}
\usepackage{enumitem}
\usepackage{microtype}

\graphicspath{{../docs/figures/}}

% ── Color palette ───────────────────────────────────────
\definecolor{yes}{HTML}{2980B9}
\definecolor{maybe}{HTML}{E67E22}
\definecolor{no}{HTML}{E74C3C}
\definecolor{epiorange}{HTML}{FF9800}
\definecolor{darktext}{HTML}{2C3E50}
\definecolor{blockbg}{HTML}{FFFFFF}

% ── Theme ───────────────────────────────────────────────
\usetheme{Default}

\definecolorstyle{SCIBColors}{
  \definecolor{colorOne}{HTML}{2C3E50}
  \definecolor{colorTwo}{HTML}{3498DB}
  \definecolor{colorThree}{HTML}{ECF0F1}
}{
  \colorlet{backgroundcolor}{colorThree}
  \colorlet{framecolor}{colorOne}
  \colorlet{titlefgcolor}{white}
  \colorlet{titlebgcolor}{colorOne}
  \colorlet{blocktitlebgcolor}{colorTwo}
  \colorlet{blocktitlefgcolor}{white}
  \colorlet{blockbodybgcolor}{blockbg}
  \colorlet{blockbodyfgcolor}{darktext}
  \colorlet{notefgcolor}{darktext}
  \colorlet{notebgcolor}{epiorange!15}
  \colorlet{notefrcolor}{epiorange}
}
\usecolorstyle{SCIBColors}

\defineblockstyle{SCIBBlock}{
  titlewidthscale=1.0, bodywidthscale=1.0, titleleft,
  titleoffsetx=0pt, titleoffsety=0pt,
  bodyoffsetx=0pt, bodyoffsety=0pt, bodyverticalshift=0pt,
  roundedcorners=12, linewidth=2pt, titleinnersep=10pt, bodyinnersep=15pt
}{
  \draw[rounded corners=\blockroundedcorners, color=blocktitlebgcolor,
        fill=blockbodybgcolor, line width=\blocklinewidth]
    (blockbody.south west) rectangle (blocktitle.north east);
  \ifBlockHasTitle
    \draw[rounded corners=\blockroundedcorners,
          fill=blocktitlebgcolor, color=blocktitlebgcolor,
          line width=\blocklinewidth]
      (blocktitle.south west) rectangle (blocktitle.north east);
  \fi
}
\useblockstyle{SCIBBlock}

\definetitlestyle{SCIBTitle}{
  width=\paperwidth, roundedcorners=0, linewidth=0pt, innersep=20pt,
  titletotopverticalspace=15mm, titletoblockverticalspace=25mm
}{
  \draw[fill=titlebgcolor, color=titlebgcolor]
    (\titleposleft, \titleposbottom)
    rectangle (\titleposright, \titlepostop);
}
\usetitlestyle{SCIBTitle}

\title{%
  \parbox{0.92\linewidth}{\centering\bfseries
    Construct Validity Failure in Single-Cell Embedding Evaluation:\\[6pt]
    Null Saturation Bounds and Source Classifier Confidence}}
\author{Elliot Tower\quad\texttt{elliot@elliottower.ai}}
\institute{}

\begin{document}

\maketitle[width=\paperwidth]

% ════════════════════════════════════════════════════════
% ROW 1
% ════════════════════════════════════════════════════════
\begin{columns}
\column{0.33}

\block{The Problem}{
  Single-cell foundation model benchmarks consistently find that
  \textbf{PCA on raw counts matches or exceeds pretrained models} on
  most evaluation axes.

  \vspace{6mm}
  Is this a genuine limitation of current models, or a construct
  validity failure in the evaluation metrics?

  \vspace{6mm}
  \textbf{Metrics serve two logically independent roles:}
  \begin{itemize}[leftmargin=*, itemsep=3mm]
    \item \textbf{Ranking}: which model embeds better?
    \item \textbf{Certification}: does a model capture biological
      signal at all?
  \end{itemize}

  \vspace{4mm}
  A metric can rank correctly while scoring every trained model
  below unstructured noise.

  \vspace{6mm}
  We evaluate \textbf{12 metrics} across \textbf{8 embedding models},
  \textbf{25 human tissues}, and \textbf{104 conditions} in CELLxGENE
  Census. Ground truth: cross-technology cell-type transfer F1 (10x
  Chromium $\to$ Smart-seq2).
}

\column{0.34}

\block{CKA Null Saturation}{
  CKA on cell-type centroids ranks embeddings well (partial
  $\rho = 0.53$). But the analytic expectation under a Gaussian null
  admits a closed form:

  \vspace{6mm}
  \begin{center}
  {\Large $\displaystyle \mathbb{E}[\mathrm{CKA}] \approx 1 -
  1.06\,\frac{k-1}{d+k}$}
  \end{center}

  \vspace{6mm}
  This null depends only on the number of cell types $k$ and
  embedding dimension $d$, independent of biology.

  \vspace{6mm}
  \begin{center}
  \begin{tabular}{@{}rcccc@{}}
    \toprule
    & \multicolumn{4}{c}{\textbf{Analytic CKA null at dimension $d$}} \\
    \cmidrule(lr){2-5}
    $k$ & $d{=}50$ & $d{=}200$ & $d{=}512$ & $d{=}1280$ \\
    \midrule
    8  & 0.872 & 0.964 & 0.986 & 0.994 \\
    14 & 0.785 & 0.936 & 0.974 & 0.989 \\
    22 & 0.691 & 0.900 & 0.958 & 0.983 \\
    44 & 0.515 & 0.813 & 0.921 & 0.965 \\
    \bottomrule
  \end{tabular}
  \end{center}

  \vspace{4mm}
  At $d \geq 512$, the null exceeds every trained-model CKA score.
  \textbf{No foundation model at native dimensionality clears the
  null in any tissue.}
}

\column{0.33}

\block{Dimensionality Is the Causal Variable}{
  The null saturation gradient is monotone across four levels of $d$:

  \vspace{6mm}
  \begin{center}
  \begin{tabular}{@{}rccc@{}}
    \toprule
    $d$ & Above null & Total & \% \\
    \midrule
    50   & 218 & 253 & \textcolor{yes}{\textbf{86\%}} \\
    200  &  76 & 456 & \textcolor{maybe}{\textbf{17\%}} \\
    512  &   0 & 125 & \textcolor{no}{\textbf{0\%}} \\
    $\geq 768$ & 0 & 4 & \textcolor{no}{\textbf{0\%}} \\
    \bottomrule
  \end{tabular}
  \end{center}

  \vspace{6mm}
  \textbf{Confound control:} scVI (foundation model, $d{=}50$) clears
  the null in 64\% of tissues. Bag-of-genes PCA ($d{=}512$) fails in
  all 25. Model family is not the driver.

  \vspace{6mm}
  \textbf{PCA-reduction:} Geneformer at native $d{=}512$ clears 0/22
  tissues. After PCA to $d{=}50$: 18/22 (82\%). Same embeddings, same
  model---dimensionality crosses the boundary.

  \vspace{6mm}
  Attempted repair ($z$-score normalization) degrades ranking by
  removing real $k$-dependent signal. The fix is the usability
  criterion, not a rescaled score.
}

\end{columns}

% ════════════════════════════════════════════════════════
% ROW 2
% ════════════════════════════════════════════════════════
\begin{columns}
\column{0.33}

\block{Relational Metrics Fail Structurally}{
  Relational consistency scores (RCS) operate on $k \times k$ cell-type
  relationship matrices.

  \vspace{6mm}
  \textbf{Advantage:} dimensionality-stable, no null saturation.

  \vspace{4mm}
  \textbf{Failure:} predict transfer only marginally ($\rho \leq 0.32$).
  All five RCS variants produce $\rho \in [0.18, 0.32]$.

  \vspace{6mm}
  The failure is structural: relational metrics measure geometric
  similarity between cell-type relationship matrices, which does not
  reflect cluster separability. Ablating all four design differences
  between RCS and official scGraph confirms no single choice explains
  the gap.

  \vspace{6mm}
  Proxy A-distance, the metric most directly grounded in
  domain-adaptation theory, fails entirely ($\rho = 0.003$, $p = 0.36$).
  Marginal distributional alignment does not determine class-conditional
  transfer.
}

\column{0.34}

\block{Source Classifier Confidence}{
  SCC trains a cell-type classifier on source embeddings, applies it
  to target cells, and reports the \textbf{mean maximum predicted
  probability}.

  \vspace{6mm}
  \begin{center}
  \begin{tabular}{@{}lccc@{}}
    \toprule
    \textbf{Metric} & $\rho$ & \textbf{95\% CI} & \textbf{Pairwise} \\
    \midrule
    SCC (LR)  & \textcolor{yes}{\textbf{+0.77}} & $[+.64,\;+.85]$ & 82.6\% \\
    SCC (RF)  & +0.59 & $[+.38,\;+.74]$ & 79.7\% \\
    SCC (SVM) & +0.52 & $[+.30,\;+.69]$ & 72.5\% \\
    SCC (kNN) & +0.51 & $[+.31,\;+.66]$ & 78.3\% \\
    \midrule
    MMD       & +0.60 & $[+.42,\;+.73]$ & 87.0\% \\
    CCAL      & +0.60 & $[+.43,\;+.72]$ & 89.9\% \\
    \midrule
    RCS (best)& +0.32 & $[+.07,\;+.52]$ & 65.2\% \\
    PAD       & +0.11 & $[-.13,\;+.34]$ & 51.6\% \\
    \bottomrule
  \end{tabular}
  \end{center}

  \vspace{5mm}
  Four classifier families all pass BH correction ($p < 0.001$),
  ruling out shared-classifier confounding.
}

\column{0.33}

\block{Convergent Validity}{
  SCC is validated against a \textbf{non-classification ground truth}:
  marker gene recovery scored in gene space via Leiden clustering.

  \vspace{6mm}
  \begin{center}
  \begin{tabular}{@{}lcc@{}}
    \toprule
    \textbf{Ground truth} & $\rho$ & \textbf{95\% CI} \\
    \midrule
    Transfer F1     & +0.77 & $[+.64,\;+.85]$ \\
    Marker gene rec.\ & +0.32 & $[+.15,\;+.44]$ \\
    \bottomrule
  \end{tabular}
  \end{center}

  \vspace{6mm}
  SCC predicts both classification-based and gene-expression-based
  ground truths, confirming that it measures an embedding-space
  property rather than a classifier artifact.

  \vspace{6mm}
  \textbf{Noise calibration:} 10 seeds $\times$ 22 tissues reveals
  that apparent non-monotonicity in Leiden-dependent metrics is
  clustering jitter. ARI non-monotonicity drops from 70\% to 6.7\%
  after permutation calibration (10.5$\times$ inflation in raw
  scores).
}

\end{columns}

% ════════════════════════════════════════════════════════
% ROW 3
% ════════════════════════════════════════════════════════
\begin{columns}
\column{0.33}

\block{The Metric Landscape}{
  \begin{center}
  \begin{tabular}{@{}lcc@{}}
    \toprule
    \textbf{Metric} & \textbf{Ranks?} & \textbf{Certifies?} \\
    \midrule
    Procrustes    & \textcolor{yes}{$\rho = 0.65$} & untested \\
    Silhouette    & \textcolor{yes}{$\rho = 0.60$} & untested \\
    CKA           & \textcolor{yes}{$\rho = 0.53$} & \textcolor{no}{saturated} \\
    NMI           & \textcolor{yes}{$\rho = 0.55$} & jitter \\
    ARI           & \textcolor{maybe}{$\rho = 0.37$} & jitter \\
    RCS           & \textcolor{no}{$\rho \leq 0.32$} & \textcolor{yes}{stable} \\
    Dir.\ stab.   & \textcolor{no}{$\rho = 0.02$} & --- \\
    \midrule
    SCC (LR)      & \textcolor{yes}{$\rho = 0.77$} & \textcolor{yes}{valid} \\
    \bottomrule
  \end{tabular}
  \end{center}

  \vspace{5mm}
  No existing metric both ranks and certifies. SCC is the strongest
  candidate for both roles.
}

\column{0.34}

\block{What This Does Not Show}{
  \textbf{Single ground truth.} Cross-technology transfer F1 is one
  downstream task. SCC's predictive power may not extend to other
  applications (trajectory inference, perturbation prediction,
  gene-program recovery).

  \vspace{5mm}
  \textbf{PBMC non-replication.} CKA's ranking ability replicates on
  pancreas and Tabula Sapiens but not PBMC, where cell-type
  classification is uniformly easy and ceiling effects compress the
  ranking signal.

  \vspace{5mm}
  \textbf{$k$ dependence.} The analytic null requires $k$ shared cell
  types across technologies. Tissues with few shared types ($k < 8$)
  were excluded.

  \vspace{5mm}
  \textbf{Causal direction.} SCC predicts transfer F1; the analysis
  does not establish whether high classifier confidence causes
  good transfer or merely correlates with it.
}

\column{0.33}

\block{Takeaway}{
  \textbf{CKA on cell-type centroids cannot certify biological signal
  at foundation-model dimensionalities.} The analytic null is a
  closed-form function of $k$ and $d$ that every current model fails
  to clear.

  \vspace{6mm}
  \begin{itemize}[leftmargin=*, itemsep=4mm]
    \item Null saturation is dimensionality-driven, not model-driven
    \item Relational metrics avoid saturation but predict transfer
      only marginally ($\rho \leq 0.32$)
    \item SCC achieves $\rho = 0.77$ across four classifier families
    \item Convergent validity from marker gene recovery ($\rho = 0.32$)
    \item Published rankings derived from relational metrics alone
      warrant re-evaluation
  \end{itemize}

  \vspace{8mm}
  \begin{center}
  \texttt{elliot@elliottower.ai}

  \vspace{3mm}
  {\small Code, data and pre-registrations:}

  \texttt{github.com/elliottower/scib-construct-validity}
  \end{center}
}

\end{columns}

\end{document}
